
où les
flèches verticales sont les homomorphismes canoniques \(K^{\bullet}\to K_\bullet\) (IV 2.5.1).

\subsection*{7.7. Spécialisation pour \(K^\bullet\) et \(\Gr^\bullet\).}
Lorsque \(X_t\) et \(X_s\) sont réguliers (par exemple si \(f\) est lisse), les deux flèches verticales précédentes sont des isomorphismes (IV 2.5), et l'homomorphisme de spécialisation peut alors s'interpréter comme un homomorphisme
\begin{equation}\tag{7.7.1}
  \sigma^\bullet:K^\bullet(X_t)\longrightarrow K^\bullet(X_s),
\end{equation}
qui peut se décrire encore comme l'unique homomorphisme rendant commutatif le diagramme suivant
\begin{equation}\tag{7.7.2}
\begin{tikzcd}[column sep=large,row sep=large]
K^\bullet(X) \arrow[r,"j^*"] \arrow[d,"i^*"'] & K^\bullet(X_t) \arrow[dl,"\sigma^\bullet"]\\
K^\bullet(X_s) & {.}
\end{tikzcd}
\end{equation}
Comme dans ce dernier \(i^*\) et \(j^*\) sont des homomorphismes d'anneaux (IV 2.7 b) avec \(j^*\) surjectif, il s'ensuit que \emph{l'homomorphisme de spécialisation} (7.7.1) \emph{est un homomorphisme d'anneaux}. Pour la même raison, \emph{si} \(X\) \emph{est séparé}, \emph{l'homomorphisme} (7.7.1) \emph{est même un homomorphisme de} \(\lambda\)-\emph{anneaux}. Pour préciser ce point, il faut seulement noter qu'un schéma noethérien régulier séparé est divisoriel (II 2.2.7.1), donc son \(K^\bullet\) s'identifie à son \(K^\bullet\) «~naïf~» (IV 2.9 (ii)), lequel est muni de façon naturelle d'une structure de \(\lambda\)-anneau (VI 3.2). D'autre part, le fait que \(X_t\) et \(X_s\) soient réguliers implique le même fait pour \(X\) (EGA \(0_{\mathrm{IV}}\) 17.1.8), et les homomorphismes \(i^*:K^\bullet(X)\to K^\bullet(X_s)\) et \(j^*:K^\bullet(X)\to K^\bullet(X_t)\) sont des \(\lambda\)-homomorphismes, donc il en est de même de (7.7.1) qui se déduit de \(i^*\) par passage au quotient.

